\section{Gaussian heat dissipation}\label{sec:heat}
We now weight the three edge matrices by Gaussians.  For the variance rates
chosen below, the sum of the star energies decreases after integration over
a plane of centers.  This decrease controls the matrix derivatives that will
produce the odd kernel.  The estimate must be
independent of the number of matrix entries; this uniformity will allow us
to pass from finite matrices to functions.

Continuous Gaussian telescoping for entangled forms appears in
Durcik~\cite[Lemma~3 and Section~3]{Durcik2015}. Here we use the Hessian of a matrix
energy and compare its costs at neighboring vertices.

Fix $h>0$, an integer $N\ge0$, and nodes $u_i=ih$, $-N\le i\le N$.
At vertex $v\in\mathbb Z/3\mathbb Z$ take $H_v=\C^{2N+1}$, and prescribe
arbitrary complex edge values $A_v(i,j)$.  For centers $p_v\in\R$ and
variances $\sigma_v>0$, define
\begin{gather*}
 g_\sigma(t)=(2\pi\sigma)^{-1/2}e^{-t^2/(2\sigma)},\qquad
 w_v(i)=h g_{\sigma_v}(u_i-p_v),\\
 W_v(i,j)=A_v(i,j)\sqrt{w_v(i)w_{v+1}(j)}.
\end{gather*}
Use the stars $R_v=[W_v\ W_{v-1}^*]$ and Grams
$T_v=R_v^*R_v$, $S_v=R_vR_v^*$ from Section~\ref{sec:trace}, and set
\[
 e_v=\Phi(T_v)=\Phi(S_v).
\]
Define the diagonal score and its insertions into the incident stars by
\begin{align*}
 L_w&=\diag_i\frac{u_i-p_w}{\sigma_w},
 &U_v&=R_v^*L_vR_v,\\
 U_{v,v+1}&=W_vL_{v+1}W_v^*,
 &U_{v,v-1}&=W_{v-1}^*L_{v-1}W_{v-1}.
\end{align*}
Define their quadratic costs and the mixed term by
\begin{equation}\label{eq:heat-costs}
 V_v=b_{T_v}(U_v),\qquad
 Y_{v,w}=b_{S_v}(U_{v,w})\quad(w\ne v),\qquad
 Z_v=\mathcal B_{S_v}(U_{v,v+1},U_{v,v-1}).
\end{equation}
Every insertion is supported on the Gram used in its definition: an
inserted Gram $A^*BA$ annihilates $\ker A$ on both sides.

For $d\in\R$, let $\Pi_d=\{p_0+p_1+p_2=d\}$, with the measure
\begin{equation}\label{eq:heat-plane-measure}
 \int_{\Pi_d}H\,\dd\pi_d
 =\int_{\R^2}H(p_0,p_1,d-p_0-p_1)\,\dd p_0\,\dd p_1.
\end{equation}
Any other pair of free coordinates gives the same measure, since the
change of coordinates has absolute determinant one.  Put $\lambda=11/10$.

We also write
\[
 \mathcal S(A)=\sum_{v=0}^2\sum_{i,j=-N}^Nh^2|A_v(i,j)|^3.
\]
Our goal in this section is the fixed-array inequality
\[
 -\partial_sJ(d,s)\ge c_*\int_{\Pi_d}\sum_vV_v\,\dd\pi_d,
 \qquad J(d,s)=\int_{\Pi_d}\sum_ve_v\,\dd\pi_d,
\]
when the variances are $\sigma_i=\alpha_i s$ and
$\boldsymbol\alpha$ is a permutation of $(1,\lambda,\lambda)$.
The odd-window estimate in Section~\ref{sec:maximal} uses wider Gaussian
weights for uniformly bounded diagonal factors, but compares each score cost
after narrowing only its own vertex. The fixed enlargement
$\lambda=11/10$ retains strict dissipation for all three comparisons.
Lemma~\ref{lem:plane-dissipation} gives an absolute $c_*>0$;
Lemma~\ref{lem:initial-energy} bounds $J$ by $C\mathcal S(A)$ at scales
$s\ge h^2$. Later, when entries change at prescribed scales, the same
differential inequality will apply between successive changes.

We first differentiate the star energies, choosing the Gram orientation
whose support stays fixed. We then integrate over the center plane and
compare neighboring stars. These two steps produce the strict decrease
that controls the score costs.

\subsection{Differentiating the star energies}

\begin{lemma}\label{lem:star-smoothness}
For fixed finite edge arrays, each $e_v$ is smooth in all centers and
positive variances.  If only the parameters at vertex $v$ vary, the support
of $T_v$ is fixed.  If the parameters at either or both other vertices
vary while those at $v$ are fixed, the support of $S_v$ is fixed.
\end{lemma}
\begin{proof}
The weighted star has the form $R_v=D_{\rm row}R_v^0D_{\rm col}$, where
$R_v^0=[A_v\ A_{v-1}^*]$ is fixed and both diagonal multipliers are strictly
positive; the column multiplier has one block for each neighbor.
Invertible row multiplication preserves the kernel, and invertible column
multiplication preserves the range.  This proves the two support
assertions, since
\[
 \ker R_v=\ker(R_v^0D_{\rm col}),\qquad
 \operatorname{ran}R_v=D_{\rm row}\operatorname{ran}R_v^0
\]
when the respective other multiplier is held fixed.

For joint smoothness, the rank-zero case is immediate.  Otherwise factor
$R_v^0=XY^*$ with $X,Y$ of full column rank.  Write
$X'=D_{\rm row}X$, $Y'=D_{\rm col}Y$ and
$P'=X'^*X'$, $Q'=Y'^*Y'$.  These last two matrices are positive definite
and depend smoothly on the parameters.  The nonzero eigenvalues of
$R_vR_v^*=X'Q'X'^*$ are those of $Q'^{1/2}P'Q'^{1/2}$, by comparing
$AA^*$ and $A^*A$ for $A=X'Q'^{1/2}$.  Therefore
\[
 e_v=\Phi(Q'^{1/2}P'Q'^{1/2}),
\]
which is smooth on the positive definite cone.
\end{proof}

\begin{lemma}\label{lem:heat-identities}
Fix a positive vector $\boldsymbol\alpha$ and set $\sigma_i=\alpha_i s$.
For all centers and all $s>0$,
\begin{align}
 2\partial_s e_v
 &=\sum_i\alpha_i\partial_{p_i}^2e_v
       -\alpha_vV_v-\sum_{w\ne v}\alpha_wY_{v,w},
       \label{eq:heat-identity}\\
 \partial_{p_{v+1}}\partial_{p_{v-1}}e_v&=Z_v.
       \label{eq:heat-mixed}
\end{align}
\end{lemma}
\begin{proof}
For $w=h g_\sigma(u-p)$ and $L=(u-p)/\sigma$,
\[
 \partial_pw=Lw,\qquad
 \partial_p^2w=(L^2-\sigma^{-1})w,\qquad
 \partial_\sigma w=\tfrac12\partial_p^2w.
\]
With the columns fixed, $T_v$ is linear in the row density, so
\[
 \partial_{p_v}T_v=U_v,\qquad
 \partial_{\sigma_v}T_v=\tfrac12\partial_{p_v}^2T_v.
\]
Its support is fixed by Lemma~\ref{lem:star-smoothness}.  The chain rule
and Lemma~\ref{lem:hessian} thus give
\[
 \partial_{\sigma_v}e_v
 =\tfrac12\partial_{p_v}^2e_v
    -\tfrac12\mathcal B_{T_v}(U_v,U_v)
 =\tfrac12\partial_{p_v}^2e_v-\tfrac12V_v.
\]
For a column vertex $w\ne v$, use instead $S_v$, which is linear in that
column density and has fixed support.  This gives
\[
 \partial_{p_w}S_v=U_{v,w},\qquad
 \partial_{\sigma_w}e_v=\tfrac12\partial_{p_w}^2e_v-\tfrac12Y_{v,w}.
\]
Summing the three variance derivatives with coefficients $\alpha_i$
proves \eqref{eq:heat-identity}.

The two summands in $S_v=W_vW_v^*+W_{v-1}^*W_{v-1}$ depend on different
column densities.  Its mixed derivative in the two column centers is
therefore zero.  Only the Hessian term remains in the mixed derivative of
$e_v$, giving \eqref{eq:heat-mixed}.
\end{proof}

\subsection{Comparing neighbors and integrating centers}

\begin{lemma}\label{lem:neighbors}
For every choice of the weighted matrices and every vertex $w$,
\begin{equation}\label{eq:heat-neighbors}
 Y_{w+1,w}+Y_{w-1,w}\le\sqrt2\,V_w.
\end{equation}
\end{lemma}
\begin{proof}
Abbreviate $A=W_w$, $B=W_{w-1}^*$, and $L=L_w$.  The inequality
$\lVert Az+Bz'\rVert^2\le2\lVert Az\rVert^2+2\lVert Bz'\rVert^2$
implies
\[
 T_w\le2Q,\qquad Q=\diag(A^*A,B^*B).
\]
Since $U_w=[A\ B]^*L[A\ B]$ is supported on $T_w$,
Lemma~\ref{lem:order} gives $b_{T_w}(U_w)\ge2^{-1/2}b_Q(U_w)$.
In a block-preserving eigenbasis of $Q$, the formula for $b_Q$ is a sum
of nonnegative squared entries.  Discarding the off-diagonal blocks yields
\[
 b_Q(U_w)\ge b_{A^*A}(A^*LA)+b_{B^*B}(B^*LB).
\]
Here $A^*A\le S_{w+1}$ and $B^*B\le S_{w-1}$, and the respective
insertions are $U_{w+1,w}$ and $U_{w-1,w}$.  Each is supported on its
smaller edge Gram.  Apply Lemma~\ref{lem:order} again, now with comparison
factor one, to obtain \eqref{eq:heat-neighbors}.
\end{proof}

\begin{lemma}\label{lem:plane-dissipation}
Fix finite edge arrays and a positive vector $\boldsymbol\alpha$, and
evaluate the energies at $\sigma_i=\alpha_i s$.  Set
\[
 J_v(d,s)=\int_{\Pi_d}e_v\,\dd\pi_d,
 \qquad J(d,s)=\sum_vJ_v(d,s),
 \qquad \Sigma_\alpha=\sum_i\alpha_i.
\]
These integrals are finite, continuously differentiable in $s>0$, and
twice continuously differentiable in $d$.  Moreover,
\begin{equation}\label{eq:heat-plane-dissipation}
 2\partial_sJ(d,s)\le\int_{\Pi_d}\sum_w
 \left[-\alpha_w+
 \sqrt2\max\{0,\Sigma_\alpha/2-\alpha_w\}\right]V_w\,\dd\pi_d.
\end{equation}
In particular, for each permutation of $(1,\lambda,\lambda)$,
\begin{equation}\label{eq:heat-strict}
 -\partial_sJ(d,s)\ge c_*\int_{\Pi_d}\sum_wV_w\,\dd\pi_d,
 \qquad
 c_* = \frac12\min\{1-\tfrac35\sqrt2,\tfrac{11}{10}-\tfrac12\sqrt2\}>0.
\end{equation}
\end{lemma}
\begin{proof}
We first justify differentiation at the fixed mesh.  On compact positive
$s$-intervals the nodes are bounded, the scores grow at most linearly in
the centers, and the first and second density ratios grow at most
quadratically.  Lemma~\ref{lem:insertion} gives
\[
 V_v\le\tfrac34\lVert L_v\rVert_{\op}^2e_v,
 \qquad Y_{v,w}\le\tfrac34\lVert L_w\rVert_{\op}^2e_v.
\]
For the second inequality, pad $L_w$ by zero on the other column block,
obtaining $B_w$ with $U_{v,w}=R_vB_wR_v^*$, and apply the lemma to
$R_v^*$.  Also $2|Z_v|\le Y_{v,v+1}+Y_{v,v-1}$.
For a Hermitian matrix $B$,
\[
 |\tr(\sqrt{T_v}R_v^*BR_v)|\le\lVert B\rVert_{\op}e_v,
\]
since $R_v\sqrt{T_v}R_v^*$ is positive with trace $e_v$.  The corresponding
column estimate follows in the same way with $S_v$ and a padded column
insertion.  Together with the density-ratio formulas in
Lemma~\ref{lem:heat-identities}, these bounds control every first center
derivative, pure second center derivative, mixed column derivative, and
first $s$-derivative by a polynomial in the centers times $e_v$.

To dominate the energy itself, use
\begin{equation}\label{eq:heat-energy-bound}
 e_v\le(\lVert W_v\rVert_{\HS}^2+
            \lVert W_{v-1}\rVert_{\HS}^2)^{3/2}
 \le\sqrt2(\lVert W_v\rVert_{\HS}^3+
                  \lVert W_{v-1}\rVert_{\HS}^3).
\end{equation}
Each edge norm cubed is bounded by a Gaussian in its two endpoint
centers, locally uniformly for $s>0$.  Indeed, the finite set of nodes is
bounded and $(u-p)^2\ge p^2/2-u^2$.  On $\Pi_d$ with $|d|\le D$, any
pair of centers satisfies
\[
 p_0^2+p_1^2+p_2^2\le4(p_i^2+p_j^2)+3D^2.
\]
Thus an endpoint-pair Gaussian, multiplied by any of the preceding
polynomials, has an integrable majorant in any fixed plane chart, locally
uniformly in $(d,s)$.  This proves all needed differentiation statements
by dominated convergence.  The constants in this justification may
depend on the finite mesh and the compact parameter intervals; they will
not enter the uniform estimate.

We can now differentiate the plane integrals. It remains to express
their second derivatives in $d$ using both pure and mixed center
derivatives, then apply Lemma~\ref{lem:neighbors} to the resulting costs.

Choosing $p_i$ as the dependent coordinate in
\eqref{eq:heat-plane-measure} gives
\[
 \partial_d^2J_v=\int_{\Pi_d}\partial_{p_i}^2e_v\,\dd\pi_d
 \qquad\text{for every }i.
\]
For the column indices $j=v+1$, $k=v-1$, first choose $p_j$ dependent to
write $\partial_dJ_v=\int_{\Pi_d}\partial_{p_j}e_v\,\dd\pi_d$.
Reparametrize this last integral with $p_k$ dependent, then differentiate
again.  It follows from \eqref{eq:heat-mixed} that
\begin{equation}\label{eq:heat-plane-mixed}
 \partial_d^2J_v
 =\int_{\Pi_d}\partial_{p_k}\partial_{p_j}e_v\,\dd\pi_d
 =\int_{\Pi_d}Z_v\,\dd\pi_d.
\end{equation}
These are identities of integrated derivatives; no pointwise equality
between pure and mixed derivatives is required.

Integrating \eqref{eq:heat-identity} and using
\eqref{eq:heat-plane-mixed}, we obtain
\[
 2\partial_sJ_v=\int_{\Pi_d}
 \left[\Sigma_\alpha Z_v-\alpha_vV_v
                   -\sum_{w\ne v}\alpha_wY_{v,w}\right]\dd\pi_d.
\]
Cauchy--Schwarz gives $2Z_v\le Y_{v,v+1}+Y_{v,v-1}$.  Sum over $v$ and
collect the two neighbor costs carrying the score at $w$.  Their common
coefficient is $\Sigma_\alpha/2-\alpha_w$.  If it is negative, discard
those terms; otherwise apply Lemma~\ref{lem:neighbors}.  This proves
\eqref{eq:heat-plane-dissipation}.

For the indicated permutations, $\Sigma_\alpha=16/5$.  The two possible
coefficients in \eqref{eq:heat-plane-dissipation} are
$-1+(3/5)\sqrt2$ and $-11/10+(1/2)\sqrt2$, both strictly negative. Dividing by two gives
\eqref{eq:heat-strict}.
\end{proof}

\subsection{The initial energy}
The previous lemma controls the accumulated costs by the energy lost along
the heat flow. We now bound the energy itself at scales no smaller than the
mesh. This estimate will also apply immediately after entries are switched
on or off, since it requires only a bound on their absolute values.

\begin{lemma}[Initial energy]\label{lem:initial-energy}
Let $h>0$, let $N\ge0$ be an integer, and let $A_v(i,j)$ be the finite complex arrays above.
For any $d\in\R$ and variances $\sigma_v\ge h^2$,
\begin{equation}\label{eq:initial-energy}
 \int_{\Pi_d}\sum_ve_v\,\dd\pi_d
 \le C_0\mathcal S(A),\qquad
 C_0=2\sqrt2\left(1+\sqrt{2/\pi}\right).
\end{equation}
The same bound holds if any entries are replaced by smaller entries in
absolute value before forming the weighted matrices.
\end{lemma}
\begin{proof}
For every center $p$ and variance $\sigma>0$,
\begin{equation}\label{eq:heat-mesh-mass}
 \sum_{i=-N}^Nh g_\sigma(u_i-p)
 \le1+\frac{2h}{\sqrt{2\pi\sigma}}.
\end{equation}
To see this, compare a grid point $u_i\ge p+h$ with the integral over its
preceding interval of length $h$, where the Gaussian is at least its
value at $u_i$. For $u_i\le p-h$, use the succeeding interval. These
intervals are disjoint, and the at most two remaining grid points each
contribute at most $h\sup g_\sigma$.

Since each $\sigma_v\ge h^2$, every density has total mass at most
$M=1+\sqrt{2/\pi}$. Weighted H\"older gives
\begin{align*}
 \lVert W_v\rVert_{\HS}^3
 &=\left(\sum_{i,j}|A_v(i,j)|^2w_v(i)w_{v+1}(j)\right)^{3/2}\\
 &\le M\sum_{i,j}|A_v(i,j)|^3w_v(i)w_{v+1}(j).
\end{align*}
Integrate with $p_v,p_{v+1}$ as free coordinates on $\Pi_d$. Each density
has integral $h$, so
\[
 \int_{\Pi_d}\lVert W_v\rVert_{\HS}^3\,\dd\pi_d
 \le M\sum_{i,j}h^2|A_v(i,j)|^3.
\]
Each edge occurs in exactly two stars. Equation~\eqref{eq:heat-energy-bound}
therefore proves \eqref{eq:initial-energy}. Replacing entries by smaller
ones in absolute value only reduces the last upper bounds, proving the
final assertion.
\end{proof}

For fixed arrays and a permutation $\boldsymbol\alpha$ of
$(1,\lambda,\lambda)$, Lemmas~\ref{lem:plane-dissipation} and
\ref{lem:initial-energy} now give, whenever $h^2\le s_0<s_1$,
\[
 \int_{s_0}^{s_1}\int_{\Pi_0}\sum_vV_v\,\dd\pi_0\,\dd s
 \le\frac{J(0,s_0)-J(0,s_1)}{c_*}
 \le\frac{C_0}{c_*}\mathcal S(A).
\]
The first inequality retains the final energy. In the next section, we
will add these inequalities over intervals with different masked arrays.
